\documentclass[11pt,letterpaper]{article}
\usepackage{shared/zoocover}

% Packages
\usepackage[utf8]{inputenc}
% geometry loaded by zoocover
\usepackage{amsmath,amssymb,amsthm}
\usepackage{graphicx}
\usepackage{hyperref}
\usepackage{cite}
\usepackage{booktabs}
\usepackage{array}
\usepackage{tikz}
\usepackage{pgfplots}
\usepackage{algorithm}
\usepackage{algorithmic}
\usepackage{enumitem}
\usepackage{mathtools}
\usepackage{xcolor}

\pgfplotsset{compat=1.18}

% Theorem environments
\newtheorem{theorem}{Theorem}
\newtheorem{lemma}[theorem]{Lemma}
\newtheorem{proposition}[theorem]{Proposition}
\newtheorem{corollary}[theorem]{Corollary}
\theoremstyle{definition}
\newtheorem{definition}{Definition}
\newtheorem{example}{Example}
\theoremstyle{remark}
\newtheorem{remark}{Remark}

% Custom commands
\newcommand{\Zoo}{\textsc{Zoo}}
\newcommand{\Hanzo}{\textsc{Hanzo}}
\newcommand{\Lux}{\textsc{Lux}}
\newcommand{\TChain}{\textsc{T-Chain}}
\newcommand{\KChain}{\textsc{K-Chain}}
\newcommand{\TFHE}{\textsc{TFHE}}
\newcommand{\CKKS}{\textsc{CKKS}}
\newcommand{\MLKEM}{\textsc{ML-KEM}}
\newcommand{\NTT}{\textsc{NTT}}
\newcommand{\LWE}{\textsc{LWE}}
\newcommand{\RLWE}{\textsc{RLWE}}

% Title information
\title{\textbf{Threshold Fully Homomorphic Encryption on Zoo Network:\\GPU-Accelerated Confidential Computation}\\
\large Privacy-Preserving Smart Contracts via TFHE and CKKS\\
\vspace{5pt}

\author{
Antje Worring, Zach Kelling\thanks{zach@zoo.ngo}\\
\textit{Zoo Labs Foundation}\\
\texttt{research@zoo.ngo}
}

\date{September 2025}

\begin{document}

\zoocoverpage

\begin{abstract}
We present \Zoo{} Network's integration of threshold fully homomorphic encryption (\TFHE{} and \CKKS{}) into a Layer~2 blockchain, enabling arbitrary computation on encrypted on-chain state. \Zoo{} deploys programmable bootstrapping (\TFHE{}) for boolean and integer circuits and approximate arithmetic (\CKKS{}) for privacy-preserving machine-learning inference directly within smart contracts. By mapping the Number Theoretic Transform butterfly network onto GPU thread blocks, we achieve a \textbf{40$\times$ speedup} over CPU-only bootstrapping, reducing single-gate bootstrap latency from 12\,ms to 0.3\,ms. Ciphertext operations are distributed across validators via an Encrypt-to-Shares (E2S) threshold decryption protocol with a \textbf{67-of-100} default quorum, ensuring that no individual validator ever possesses a complete decryption key. Key shares are protected in transit with \MLKEM{}-768 post-quantum key encapsulation and anchored on \Lux{} \KChain{}. Formal correctness of the \CKKS{} rescaling pipeline, \TFHE{} noise budget tracking, and GPU scaling bounds are verified in Lean~4 (proofs reside in \texttt{Crypto/FHE/CKKS.lean}, \texttt{Crypto/FHE/TFHE.lean}, and \texttt{GPU/FHEScaling.lean}). We benchmark confidential DeFi order matching, sealed-bid auctions, and encrypted neural-network inference, demonstrating practical throughput for production deployment on \Zoo{} \TChain{}.
\end{abstract}

\tableofcontents
\newpage

%----------------------------------------------------------------------
\section{Introduction}
\label{sec:intro}
%----------------------------------------------------------------------

Public blockchains expose every transaction and every byte of contract state to all participants. This transparency, while valuable for auditability, precludes entire classes of applications: sealed-bid auctions, private order matching, confidential voting, and machine-learning inference over sensitive data.

Two cryptographic families address this gap:

\begin{enumerate}
  \item \textbf{Zero-knowledge proofs} (ZKPs) allow a prover to convince a verifier of a statement's truth without revealing witness data. ZKPs prove facts about plaintext but do not enable \emph{computation} on encrypted data.
  \item \textbf{Fully homomorphic encryption} (FHE) permits arbitrary computation directly on ciphertexts. The result, once decrypted, equals the output of the same computation on plaintexts~\cite{gentry2009fhe}.
\end{enumerate}

FHE is strictly more expressive than ZKPs for on-chain privacy: a smart contract can evaluate functions over encrypted inputs without any party---including validators---ever observing plaintext values. \Zoo{} Network combines two complementary FHE schemes:

\begin{itemize}
  \item \textbf{\TFHE{}} (Torus FHE): bit-level and small-integer operations with fast programmable bootstrapping, ideal for branching logic, comparisons, and access-control circuits.
  \item \textbf{\CKKS{}} (Cheon--Kim--Kim--Song): approximate fixed-point arithmetic over vectors of real numbers, ideal for neural-network inference, statistical aggregation, and financial modelling.
\end{itemize}

\Zoo{}'s approach differs from prior work in three ways: (1)~GPU-accelerated bootstrapping that reduces latency by 40$\times$, making FHE practical inside block times; (2)~threshold decryption integrated with \Zoo{} \TChain{} validators so that decryption requires 67-of-100 consensus; and (3)~formal verification in Lean~4 of noise budgets, rescaling correctness, and GPU parallelism bounds.

This paper is organized as follows. Sections~\ref{sec:tfhe}--\ref{sec:ckks} describe the two FHE schemes and their on-chain encoding. Section~\ref{sec:gpu} details GPU acceleration. Section~\ref{sec:threshold} presents the threshold decryption protocol. Section~\ref{sec:tchain} explains integration with \TChain{}. Section~\ref{sec:formal} covers formal verification. Section~\ref{sec:apps} surveys applications, Section~\ref{sec:bench} provides benchmarks, and Section~\ref{sec:related} discusses related work.

%----------------------------------------------------------------------
\section{TFHE on Zoo}
\label{sec:tfhe}
%----------------------------------------------------------------------

\subsection{LWE Encryption}

\TFHE{} operates over the Learning With Errors (\LWE{}) problem. A plaintext bit $m \in \{0,1\}$ is encrypted as:

\begin{equation}
  \mathsf{ct} = (\mathbf{a}, b) \in \mathbb{Z}_q^n \times \mathbb{Z}_q, \quad b = \langle \mathbf{a}, \mathbf{s} \rangle + m \cdot \lfloor q/2 \rfloor + e
\end{equation}

\noindent where $\mathbf{s} \in \mathbb{Z}_q^n$ is the secret key, $\mathbf{a} \xleftarrow{\$} \mathbb{Z}_q^n$ is uniform, and $e$ is a small Gaussian noise term with standard deviation $\sigma$. \Zoo{} fixes $n = 630$, $q = 2^{32}$, and $\sigma = 2^{-15}$, yielding 128-bit security under standard \LWE{} hardness assumptions~\cite{albrecht2015lwe}.

\subsection{Boolean Gate Evaluation}

Homomorphic boolean gates (AND, OR, XOR, NOT) on \LWE{} ciphertexts follow the gate-bootstrapping paradigm of Chillotti et al.~\cite{chillotti2020tfhe}. For two input ciphertexts $\mathsf{ct}_0, \mathsf{ct}_1$:

\begin{equation}
  \mathsf{ct}_{\mathrm{AND}} = \mathsf{Bootstrap}\!\left( (0, \lfloor q/8 \rfloor) - \mathsf{ct}_0 - \mathsf{ct}_1 \right)
\end{equation}

The bootstrap operation simultaneously evaluates a lookup table and refreshes noise, producing a ciphertext with bounded noise regardless of input noise level.

\subsection{Programmable Bootstrapping}

\Zoo{} uses programmable bootstrapping (PBS) to evaluate arbitrary functions $f: \mathbb{Z}_p \to \mathbb{Z}_p$ during the bootstrap step itself. The test polynomial $v(X) \in \mathbb{Z}_q[X]/(X^N+1)$ encodes $f$ in its coefficients:

\begin{equation}
  v(X) = \sum_{i=0}^{N-1} f\!\left(\left\lfloor \frac{i \cdot p}{2N} \right\rfloor\right) X^i
\end{equation}

This eliminates post-bootstrap computation for common operations (sign extraction, ReLU, step functions), reducing circuit depth and noise accumulation.

\subsection{Noise Management}

Each homomorphic operation adds noise to the ciphertext. \Zoo{} tracks noise budgets per ciphertext in the VM state:

\begin{definition}[Noise Budget]
For ciphertext $\mathsf{ct} = (\mathbf{a}, b)$ with secret $\mathbf{s}$, the noise $e = b - \langle \mathbf{a}, \mathbf{s} \rangle - m \cdot \lfloor q/2 \rfloor$. The noise budget is $\beta = \lfloor \log_2(q/2) \rfloor - \lceil \log_2 |e| \rceil$ bits. Decryption succeeds when $\beta > 0$.
\end{definition}

The VM rejects execution when a ciphertext's projected noise budget would drop below 1 bit, forcing the contract to insert a bootstrap before proceeding. Gas metering accounts for bootstrap cost.

%----------------------------------------------------------------------
\section{CKKS on Zoo}
\label{sec:ckks}
%----------------------------------------------------------------------

\subsection{Approximate Arithmetic}

\CKKS{}~\cite{cheon2017ckks} encodes a vector $\mathbf{z} \in \mathbb{C}^{N/2}$ into a polynomial ring $R_q = \mathbb{Z}_q[X]/(X^N+1)$ via the canonical embedding inverse, scaled by $\Delta$:

\begin{equation}
  m(X) = \lfloor \Delta \cdot \sigma^{-1}(\mathbf{z}) \rceil \in R_q
\end{equation}

\noindent where $\Delta = 2^{40}$ is the scaling factor and $\sigma^{-1}$ is the inverse canonical embedding. Encryption produces $\mathsf{ct} = (\mathsf{ct}_0, \mathsf{ct}_1) \in R_q^2$ under \RLWE{}.

\subsection{Homomorphic Operations}

Addition of two \CKKS{} ciphertexts at the same level is component-wise. Multiplication produces a degree-2 ciphertext that must be relinearized:

\begin{align}
  \mathsf{ct}_{\mathrm{mult}} &= (\mathsf{ct}_0^{(a)} \cdot \mathsf{ct}_0^{(b)},\; \mathsf{ct}_0^{(a)} \cdot \mathsf{ct}_1^{(b)} + \mathsf{ct}_1^{(a)} \cdot \mathsf{ct}_0^{(b)},\; \mathsf{ct}_1^{(a)} \cdot \mathsf{ct}_1^{(b)}) \\
  \mathsf{ct}_{\mathrm{relin}} &= \mathsf{Relin}(\mathsf{ct}_{\mathrm{mult}}, \mathsf{evk})
\end{align}

\subsection{Rescaling and Level Management}

After multiplication, the scaling factor doubles to $\Delta^2$. Rescaling divides by $\Delta$ and drops one prime from the ciphertext modulus chain:

\begin{equation}
  \mathsf{Rescale}(\mathsf{ct}, q_\ell) \to \mathsf{ct}' \text{ at level } \ell - 1, \quad q_{\ell-1} = q_\ell / p_\ell
\end{equation}

\Zoo{} provisions a modulus chain of $L = 15$ levels by default, supporting up to 14 sequential multiplications---sufficient for a 4-layer neural network with ReLU approximations. The VM tracks levels alongside gas, rejecting operations that would exhaust the chain.

\subsection{ML Inference on Encrypted Data}

A neural-network forward pass on encrypted inputs uses \CKKS{} slot packing. For a weight matrix $W \in \mathbb{R}^{m \times n}$ and encrypted input $\mathsf{Enc}(\mathbf{x})$:

\begin{equation}
  \mathsf{Enc}(W\mathbf{x} + \mathbf{b}) = W \otimes \mathsf{Enc}(\mathbf{x}) \oplus \mathsf{Enc}(\mathbf{b})
\end{equation}

\noindent where $\otimes$ denotes plaintext--ciphertext multiplication and $\oplus$ denotes ciphertext addition. Activation functions are approximated by low-degree polynomials (degree-3 Chebyshev for ReLU) consuming one multiplicative level each.

%----------------------------------------------------------------------
\section{GPU Acceleration}
\label{sec:gpu}
%----------------------------------------------------------------------

\subsection{NTT Butterfly Parallelism}

The computational bottleneck in both \TFHE{} and \CKKS{} is polynomial multiplication in $R_q$, implemented via the Number Theoretic Transform (\NTT{}). The \NTT{} butterfly network of depth $\log_2 N$ maps naturally to GPU thread blocks:

\begin{itemize}
  \item Each butterfly stage is a parallel map over $N/2$ independent pairs.
  \item Stages synchronize via \texttt{\_\_syncthreads()} within a block or global barrier across blocks.
  \item For $N = 2^{14} = 16384$ (Zoo's default polynomial degree), a single SM processes 8192 butterflies per stage.
\end{itemize}

\subsection{TFHE Bootstrap on GPU}

A single \TFHE{} bootstrap consists of:

\begin{enumerate}
  \item \textbf{Blind rotation}: $n = 630$ external products, each requiring two \NTT{} forward/inverse transforms on degree-$N$ polynomials.
  \item \textbf{Key switching}: one matrix--vector product in $\mathbb{Z}_q^{n' \times n}$.
  \item \textbf{Modulus switching}: rounding from $q$ to $q'$.
\end{enumerate}

On CPU (single-threaded, AVX-512), one bootstrap takes $\sim$12\,ms. On an NVIDIA H100 GPU:

\begin{itemize}
  \item Blind rotation: 630 external products execute as a batched \NTT{} kernel, 0.22\,ms.
  \item Key switching: cuBLAS GEMV, 0.05\,ms.
  \item Modulus switching: element-wise kernel, 0.03\,ms.
  \item \textbf{Total: 0.3\,ms} (40$\times$ speedup).
\end{itemize}

\subsection{Batch Amortization}

When multiple independent bootstraps are required (e.g., evaluating a 256-bit comparison circuit), \Zoo{} batches them into a single GPU dispatch:

\begin{equation}
  T_{\mathrm{batch}}(k) = T_{\mathrm{single}} + \alpha \cdot k, \quad \alpha \approx 0.015\,\text{ms per additional bootstrap}
\end{equation}

\noindent for $k \leq 4096$ bootstraps. Beyond this, L2 cache pressure causes throughput degradation. At $k = 256$ (a typical 8-bit integer comparison), total latency is $\sim$4.1\,ms.

\subsection{Memory Bounds}

FHE ciphertexts are large. A single \CKKS{} ciphertext at level $\ell = 15$ with $N = 16384$ occupies:

\begin{equation}
  |\mathsf{ct}| = 2 \cdot N \cdot \ell \cdot 8 = 2 \cdot 16384 \cdot 15 \cdot 8 \approx 3.75\,\text{MB}
\end{equation}

An H100 with 80\,GB HBM3 can hold $\sim$21{,}000 such ciphertexts simultaneously. \Zoo{} validators must provision at least 40\,GB GPU memory for FHE workloads. The VM enforces per-contract ciphertext storage limits via gas metering.

%----------------------------------------------------------------------
\section{Threshold Decryption}
\label{sec:threshold}
%----------------------------------------------------------------------

\subsection{Motivation}

In a blockchain setting, no single entity should hold the decryption key. \Zoo{} distributes the \TFHE{}/\CKKS{} secret key across validators using Shamir secret sharing with a 67-of-100 reconstruction threshold.

\subsection{Encrypt-to-Shares Protocol}

Key generation proceeds as follows:

\begin{enumerate}
  \item A trusted dealer (the genesis ceremony) generates $\mathbf{s} \in \mathbb{Z}_q^n$ and computes Shamir shares $\mathbf{s}_i$ for $i \in \{1, \ldots, 100\}$ with threshold $t = 67$.
  \item Each share $\mathbf{s}_i$ is encapsulated under validator $i$'s \MLKEM{}-768 public key and published to \Lux{} \KChain{} as a sealed transaction.
  \item Validators decrypt their share locally and store it in a hardware security module (HSM) or TEE enclave.
\end{enumerate}

\begin{definition}[E2S Decryption]
To decrypt ciphertext $\mathsf{ct} = (\mathbf{a}, b)$, each participating validator $i$ computes a partial decryption:
\begin{equation}
  d_i = \langle \mathbf{a}, \mathbf{s}_i \rangle + e_i
\end{equation}
where $e_i$ is a smudging noise term with $\|e_i\| \ll q/4$ to prevent share leakage. The aggregator reconstructs:
\begin{equation}
  m = b - \sum_{i \in S} \lambda_i \cdot d_i \pmod{q}, \quad |S| \geq 67
\end{equation}
where $\lambda_i$ are Lagrange interpolation coefficients for the set $S$.
\end{definition}

\subsection{67-of-100 Default Threshold}

The threshold $t = 67$ is chosen to satisfy:

\begin{itemize}
  \item \textbf{Safety}: an adversary controlling fewer than 34 validators cannot decrypt any ciphertext.
  \item \textbf{Liveness}: the system tolerates up to 33 offline or Byzantine validators while maintaining decryption capability.
  \item \textbf{Alignment with BFT}: the $2/3 + 1$ threshold matches the Byzantine fault tolerance bound used in \Zoo{} \TChain{} consensus.
\end{itemize}

\subsection{Key Share Rotation}

Validator set changes trigger proactive secret sharing (PSS). New shares are computed via a resharing protocol without reconstructing $\mathbf{s}$:

\begin{equation}
  \mathbf{s}_i^{(\mathrm{new})} = \sum_{j \in S_{\mathrm{old}}} \lambda_j \cdot \mathbf{s}_{j \to i}
\end{equation}

\noindent where $\mathbf{s}_{j \to i}$ is a re-share from old validator $j$ to new validator $i$, again encapsulated under \MLKEM{}-768. This ensures forward secrecy: compromised old shares cannot decrypt ciphertexts encrypted under the new epoch's public key.

%----------------------------------------------------------------------
\section{Integration with T-Chain}
\label{sec:tchain}
%----------------------------------------------------------------------

\Zoo{} \TChain{} (Threshold Chain) combines threshold signing~\cite{zoo-threshold-signatures} with FHE in a unified VM execution environment.

\subsection{FHE Opcodes}

The \TChain{} VM extends the EVM with FHE-specific opcodes:

\begin{table}[h]
\centering
\small
\begin{tabular}{lll}
\toprule
\textbf{Opcode} & \textbf{Operation} & \textbf{Gas Cost} \\
\midrule
\texttt{FHE\_ENC} & Encrypt plaintext to \TFHE{}/\CKKS{} ciphertext & 50{,}000 \\
\texttt{FHE\_ADD} & Homomorphic addition & 10{,}000 \\
\texttt{FHE\_MUL} & Homomorphic multiplication (+ relin) & 200{,}000 \\
\texttt{FHE\_BOOT} & Programmable bootstrap (\TFHE{}) & 500{,}000 \\
\texttt{FHE\_RESCALE} & Rescale \CKKS{} ciphertext & 80{,}000 \\
\texttt{FHE\_ROTATE} & Rotate \CKKS{} slots & 150{,}000 \\
\texttt{FHE\_DEC\_REQ} & Request threshold decryption & 1{,}000{,}000 \\
\bottomrule
\end{tabular}
\caption{FHE opcodes added to the \TChain{} VM}
\label{tab:opcodes}
\end{table}

Gas costs reflect computational load. \texttt{FHE\_DEC\_REQ} is the most expensive because it triggers a multi-round threshold decryption protocol across 67+ validators.

\subsection{GPU Acceleration Configuration}

Validators configure GPU FHE acceleration in their node configuration:

\begin{itemize}
  \item \texttt{fhe.gpu.enabled}: boolean, default \texttt{true} for \TChain{} validators.
  \item \texttt{fhe.gpu.device}: CUDA device index (supports multi-GPU).
  \item \texttt{fhe.gpu.max\_batch}: maximum bootstrap batch size (default 4096).
  \item \texttt{fhe.gpu.memory\_limit}: GPU memory reserved for FHE (default 40\,GB).
\end{itemize}

Validators without GPU support may still participate in consensus and threshold decryption but cannot execute FHE opcodes. Block producers for FHE-heavy blocks are selected preferentially from the GPU-enabled validator subset.

\subsection{Ciphertext Storage}

\TFHE{} ciphertexts (LWE, $\sim$5\,KB each) and \CKKS{} ciphertexts ($\sim$3.75\,MB at full level) are stored in a dedicated ciphertext trie separate from the account state trie. Content-addressing (Blake3 hash) enables deduplication. State rent is charged per ciphertext per epoch to bound storage growth.

%----------------------------------------------------------------------
\section{Formal Verification}
\label{sec:formal}
%----------------------------------------------------------------------

\Zoo{} maintains Lean~4 proofs for critical FHE subsystems, following the broader \Zoo{} formal verification initiative~\cite{zoo-network-architecture}.

\subsection{CKKS Correctness}

The file \texttt{Crypto/FHE/CKKS.lean} proves:

\begin{theorem}[Rescaling Precision]
For a \CKKS{} ciphertext at level $\ell$ with scaling factor $\Delta$ and modulus $q_\ell$, rescaling to level $\ell - 1$ introduces approximation error bounded by:
\begin{equation}
  \|\epsilon_{\mathrm{rescale}}\|_\infty \leq \frac{N \cdot B_{\mathrm{key}}}{p_\ell}
\end{equation}
where $B_{\mathrm{key}}$ bounds the secret key coefficients and $p_\ell$ is the dropped prime.
\end{theorem}

\subsection{TFHE Noise Tracking}

The file \texttt{Crypto/FHE/TFHE.lean} verifies:

\begin{theorem}[Bootstrap Noise Bound]
After programmable bootstrapping with parameters $(n, N, q, \sigma, B_g, d_g)$, the output noise satisfies:
\begin{equation}
  \|e_{\mathrm{out}}\| \leq n \cdot d_g \cdot (N+1) \cdot B_g \cdot \sigma_{\mathrm{bsk}} + (1 + N \cdot \|\mathbf{s}\|_1) \cdot \frac{B_g^{d_g}}{2}
\end{equation}
This bound is strictly less than $q/4$ for \Zoo{}'s parameter choices, guaranteeing correct decryption.
\end{theorem}

\subsection{GPU Scaling}

The file \texttt{GPU/FHEScaling.lean} proves:

\begin{proposition}[NTT Parallelism]
The $\log_2 N$-stage butterfly network for degree-$N$ NTT admits a work-depth decomposition with:
\begin{equation}
  W = \frac{N}{2} \log_2 N, \quad D = \log_2 N
\end{equation}
Given $P \geq N/2$ parallel processors, the time complexity is $\Theta(\log_2 N)$, achieving linear speedup up to $P = N/2$.
\end{proposition}

%----------------------------------------------------------------------
\section{Applications}
\label{sec:apps}
%----------------------------------------------------------------------

\subsection{Confidential DeFi}

\textbf{Encrypted order matching.} Users submit encrypted limit orders $\mathsf{Enc}(\mathit{price}, \mathit{quantity})$. The matching engine evaluates comparison circuits homomorphically (\TFHE{} for price comparison, \CKKS{} for quantity arithmetic). Matched orders are decrypted via threshold decryption; unmatched orders remain encrypted and invisible to other participants.

This eliminates front-running, sandwich attacks, and information leakage from the order book.

\subsection{Private AI Inference}

A model owner publishes encrypted weights $\mathsf{Enc}(W_i)$ and a user submits encrypted input $\mathsf{Enc}(\mathbf{x})$. The contract evaluates the forward pass entirely under \CKKS{} encryption:

\begin{equation}
  \mathsf{Enc}(\hat{y}) = f_{\mathrm{NN}}(\mathsf{Enc}(W_1), \ldots, \mathsf{Enc}(W_L), \mathsf{Enc}(\mathbf{x}))
\end{equation}

Neither the model weights nor the input are revealed to any party. The user decrypts $\hat{y}$ locally. This enables monetization of proprietary models without exposing intellectual property.

\subsection{Sealed-Bid Auctions}

Bidders submit $\mathsf{Enc}(\mathit{bid}_i)$ to an auction contract. The contract computes the maximum via a homomorphic comparison tree (\TFHE{} with $\lceil \log_2 k \rceil$ rounds for $k$ bids). Only the winning bid is decrypted; losing bids remain sealed.

%----------------------------------------------------------------------
\section{Benchmarks}
\label{sec:bench}
%----------------------------------------------------------------------

All benchmarks run on \Zoo{} testnet validators equipped with NVIDIA H100 80\,GB GPUs, AMD EPYC 9654 CPUs (96 cores), and 512\,GB DDR5.

\subsection{Bootstrap Latency}

\begin{table}[h]
\centering
\begin{tabular}{lrrr}
\toprule
\textbf{Operation} & \textbf{CPU (ms)} & \textbf{GPU (ms)} & \textbf{Speedup} \\
\midrule
Single TFHE bootstrap & 12.0 & 0.30 & 40$\times$ \\
8-bit integer comparison & 96.0 & 4.10 & 23$\times$ \\
256-bit comparison circuit & 3{,}072 & 52.0 & 59$\times$ \\
Batch 4096 bootstraps & 49{,}152 & 62.0 & 793$\times$ \\
\bottomrule
\end{tabular}
\caption{TFHE bootstrap latency: CPU vs GPU}
\label{tab:bootstrap}
\end{table}

Batch amortization yields superlinear speedup because GPU occupancy increases with batch size until the 4096-element sweet spot.

\subsection{CKKS Multiplication Throughput}

\begin{table}[h]
\centering
\begin{tabular}{lrr}
\toprule
\textbf{Metric} & \textbf{CPU} & \textbf{GPU} \\
\midrule
Single ciphertext multiply + relin (ms) & 8.5 & 0.6 \\
Throughput (multiplications/sec) & 118 & 1{,}667 \\
Rescale (ms) & 2.1 & 0.15 \\
Slot rotation (ms) & 6.2 & 0.45 \\
\bottomrule
\end{tabular}
\caption{CKKS operation throughput at $N = 16384$, $L = 15$}
\label{tab:ckks}
\end{table}

\subsection{Threshold Decryption Latency}

\begin{table}[h]
\centering
\begin{tabular}{lr}
\toprule
\textbf{Phase} & \textbf{Latency (ms)} \\
\midrule
Partial decryption (per validator) & 0.8 \\
Network collection (67 shares, p95) & 180 \\
Lagrange aggregation & 2.1 \\
Smudging noise verification & 0.3 \\
\textbf{Total end-to-end} & \textbf{$\sim$183} \\
\bottomrule
\end{tabular}
\caption{Threshold decryption latency for LWE ciphertext (67-of-100)}
\label{tab:threshold}
\end{table}

The dominant cost is network latency for collecting 67 partial decryptions. Computation is negligible relative to communication.

\subsection{End-to-End Application Benchmarks}

\begin{table}[h]
\centering
\begin{tabular}{lrr}
\toprule
\textbf{Application} & \textbf{Latency (ms)} & \textbf{Gas Cost} \\
\midrule
Encrypted order match (2 orders) & 48 & 1.8M \\
Sealed-bid auction (64 bids) & 312 & 42M \\
4-layer NN inference (CKKS) & 580 & 68M \\
\bottomrule
\end{tabular}
\caption{End-to-end application benchmarks on GPU-enabled validators}
\label{tab:e2e}
\end{table}

%----------------------------------------------------------------------
\section{Related Work}
\label{sec:related}
%----------------------------------------------------------------------

\textbf{Lux TFHE.} The \Lux{} \texttt{lux-tfhe} paper~\cite{lux-tfhe} introduced GPU-accelerated TFHE for the Lux L0 base layer. \Zoo{} extends this work with threshold decryption, \CKKS{} support, and deeper VM integration via FHE-specific opcodes on \TChain{}.

\textbf{Zama fhEVM.} Zama's fhEVM~\cite{zama-fhevm} pioneered FHE smart contracts on EVM-compatible chains. \Zoo{} differs in three respects: (1)~\CKKS{} support for approximate arithmetic, enabling ML inference; (2)~threshold decryption rather than single-operator key management; (3)~formal verification of noise bounds and GPU scaling in Lean~4.

\textbf{Sunscreen FHE.} The Sunscreen compiler~\cite{sunscreen2023} provides a Rust-based FHE toolkit targeting BFV and CKKS. \Zoo{} builds on similar compiler infrastructure but integrates it directly into the blockchain VM with gas metering and threshold key management.

\textbf{Concrete/TFHE-rs.} Zama's \texttt{tfhe-rs} library~\cite{zama-tfhers} provides a Rust implementation of TFHE with GPU support. \Zoo{}'s GPU kernels build on this foundation, extending it with batched bootstrapping optimizations and \TChain{} integration.

\textbf{Threshold FHE in MPC.} Boneh et al.~\cite{boneh2018threshold} formalized threshold FHE for multi-party computation. \Zoo{} adapts this framework to the blockchain setting, replacing interactive MPC rounds with asynchronous partial-decryption collection anchored by BFT consensus.

%----------------------------------------------------------------------
\section{Conclusion}
\label{sec:conclusion}
%----------------------------------------------------------------------

We have presented \Zoo{} Network's threshold FHE system, combining \TFHE{} for boolean circuits and \CKKS{} for approximate arithmetic with GPU acceleration and validator-distributed threshold decryption. Key results:

\begin{enumerate}
  \item \textbf{40$\times$ GPU speedup} for single \TFHE{} bootstrap (12\,ms $\to$ 0.3\,ms), with up to 793$\times$ for batched operations.
  \item \textbf{67-of-100 threshold decryption} aligned with BFT consensus, with end-to-end latency of $\sim$183\,ms dominated by network communication.
  \item \textbf{Formal verification} of \CKKS{} rescaling precision, \TFHE{} noise bounds, and GPU parallelism in Lean~4.
  \item \textbf{Practical applications} demonstrated: encrypted order matching (48\,ms), sealed-bid auctions (312\,ms), and private neural-network inference (580\,ms).
\end{enumerate}

FHE on \Zoo{} \TChain{} enables a new class of privacy-preserving smart contracts where computation proceeds on encrypted state without any party---including validators---observing plaintext values. Combined with post-quantum key encapsulation (\MLKEM{}-768) for share distribution, the system provides long-term confidentiality against both classical and quantum adversaries.

Future work includes bootstrapping-free FHE schemes for reduced latency, hardware-accelerated key switching via custom FPGA pipelines, and integration with \Zoo{}'s experience ledger~\cite{zoo-tokenomics} for privacy-preserving model training.

\section*{Acknowledgments}

The authors thank the \Zoo{} Labs Foundation engineering team for GPU kernel development, the \Lux{} cryptography team for the \texttt{lux-tfhe} foundation, and the \Hanzo{} infrastructure team for validator provisioning.

\begin{thebibliography}{99}

\bibitem{gentry2009fhe}
Craig Gentry.
\textit{Fully Homomorphic Encryption Using Ideal Lattices}.
STOC, 2009.

\bibitem{chillotti2020tfhe}
Ilaria Chillotti, Nicolas Gama, Mariya Georgieva, and Malika Izabach\`{e}ne.
\textit{TFHE: Fast Fully Homomorphic Encryption over the Torus}.
Journal of Cryptology, 33(1):34--91, 2020.

\bibitem{cheon2017ckks}
Jung Hee Cheon, Andrey Kim, Miran Kim, and Yongsoo Song.
\textit{Homomorphic Encryption for Arithmetic of Approximate Numbers}.
ASIACRYPT, 2017.

\bibitem{albrecht2015lwe}
Martin R. Albrecht, Rachel Player, and Sam Scott.
\textit{On the Concrete Hardness of Learning with Errors}.
Journal of Mathematical Cryptology, 9(3):169--203, 2015.

\bibitem{boneh2018threshold}
Dan Boneh, Rosario Gennaro, Steven Goldfeder, Aayush Jain, Sam Kim, Peter M.R. Rasmussen, and Amit Sahai.
\textit{Threshold Cryptosystems from Threshold Fully Homomorphic Encryption}.
CRYPTO, 2018.

\bibitem{zama-fhevm}
Zama.
\textit{fhEVM: Confidential Smart Contracts on the EVM Using Fully Homomorphic Encryption}.
Zama Whitepaper, 2023.

\bibitem{zama-tfhers}
Zama.
\textit{TFHE-rs: A Pure Rust Implementation of the TFHE Scheme}.
\url{https://github.com/zama-ai/tfhe-rs}, 2023.

\bibitem{sunscreen2023}
Sunscreen.
\textit{Sunscreen: A Compiler for FHE Programs}.
Sunscreen Tech Report, 2023.

\bibitem{lux-tfhe}
Lux Industries.
\textit{GPU-Accelerated TFHE for Lux Network}.
Lux Research, 2025.

\bibitem{zoo-tokenomics}
Zach Kelling.
\textit{Zoo Network Tokenomics: Fair Distribution Through Complete Airdrop}.
Zoo Labs Foundation, 2025.

\bibitem{zoo-gpu-evm}
Zach Kelling.
\textit{Zoo EVM: A Post-Quantum Safe Layer 2 Blockchain with AI-Native Precompiles}.
Zoo Labs Foundation, 2026.

\bibitem{zoo-network-architecture}
Zach Kelling.
\textit{Zoo Network: Decentralized AI Training and Inference Layer with Semantic Optimization}.
Zoo Labs Foundation, 2025.

\bibitem{zoo-threshold-signatures}
Zach Kelling.
\textit{Threshold Signatures on Zoo T-Chain} (forthcoming).
Zoo Labs Foundation, 2026.

\end{thebibliography}

\end{document}
